% MSAE-AD Framework Overview Figure
% Uses TikZ for a professional architecture diagram
%
% This figure is designed to span two columns (figure*)
% It shows the two-phase MSAE-AD pipeline:
%   Phase 1 (Fitting): Normal images → DINOv2 → MSAE training
%   Phase 2 (Inference): Test image → DINOv2 → MSAE → Anomaly map

\begin{figure*}[t]
  \centering
  \begin{tikzpicture}[
    node distance=0.8cm,
    block/.style={
      rectangle, draw=black, thick, rounded corners=3pt,
      minimum width=1.6cm, minimum height=0.7cm,
      align=center, font=\small\sffamily, inner sep=4pt
    },
    bigblock/.style={
      rectangle, draw=black, thick, rounded corners=4pt,
      minimum width=2.0cm, minimum height=0.8cm,
      align=center, font=\small\sffamily\bfseries, inner sep=5pt,
      fill=blue!5
    },
    modelblock/.style={
      rectangle, draw=cvprblue, very thick, rounded corners=4pt,
      minimum width=2.2cm, minimum height=0.9cm,
      align=center, font=\small\sffamily\bfseries, inner sep=5pt,
      fill=cvprblue!8
    },
    smallbox/.style={
      rectangle, draw=black, thick, rounded corners=2pt,
      minimum width=1.3cm, minimum height=0.55cm,
      align=center, font=\footnotesize\sffamily, inner sep=3pt,
      fill=green!5
    },
    arr/.style={->, >=stealth, thick, draw=black!70},
    arrblue/.style={->, >=stealth, very thick, draw=cvprblue},
    dashedbox/.style={
      rectangle, draw=gray, dashed, rounded corners=6pt,
      inner sep=8pt, fill=gray!3
    },
    label/.style={font=\footnotesize\sffamily\bfseries, align=center},
    note/.style={font=\footnotesize\sffamily, align=center, text=black!60},
  ]

  % ========================================================================
  % Define coordinates for layout
  % ========================================================================
  % We split into two main columns: FIT (left) and INFERENCE (right)
  % With a dashed box around each phase

  % ---- PHASE 1: FITTING (Left Column) ----
  \node[label, text=cvprblue] (fitlabel) at (0, 2.0) {Phase 1: Fitting (once per category)};

  \node[draw=blue!30, dashed, rounded corners=8pt, fill=blue!2,
        minimum width=4.8cm, minimum height=7.2cm] (fitbox) at (0, -1.8) {};

  % Normal images (input)
  \node[block, fill=orange!10] (normalimgs) at (0, 1.0) {$k$ Normal\\Images};

  % DINOv2 feature extractor
  \node[block, fill=yellow!10] (dinov2_fit) at (0, -0.5) {Frozen\\DINOv2-G};

  % Layer aggregation
  \node[smallbox, fill=yellow!5] (layers) at (0, -1.6) {Layers 22--28\\Mean Pooling};

  % Normal patch features
  \node[block, fill=green!8] (normalfeats) at (0, -2.7) {Normal Patch\\Features\\$\mathcal{X}_{\text{normal}}$};

  % Data augmentation
  \node[note] at (1.6, -0.1) {$\times N_a$ rot.};

  % MSAE Model (the core)
  \node[modelblock, minimum width=2.8cm, minimum height=1.2cm] (msae_fit) at (0, -4.5) {\textbf{MSAE Training}\\[2pt]
    \footnotesize$\mathcal{L} = \sum_{m}\alpha_m \mathcal{L}_{\text{recon}}^m$};

  % MSAE internal - granularities
  \node[smallbox, fill=blue!5, minimum width=2.8cm] (gran) at (0, -5.8) {\footnotesize Granularities: $\{k_1, k_2, ..., k_M\}$\\[-1pt]
    \footnotesize UW / RW weighting};

  % Fitted model
  \node[modelblock, minimum width=2.4cm] (savedmodel) at (0, -6.8) {Trained MSAE\\$\{E^*, G^*\}$};

  % Arrows for fitting phase
  \draw[arr] (normalimgs) -- (dinov2_fit);
  \draw[arr] (dinov2_fit) -- (layers);
  \draw[arr] (layers) -- (normalfeats);
  \draw[arr] (normalfeats) -- (msae_fit);
  \draw[arr] (msae_fit) -- (gran);
  \draw[arr] (gran) -- (savedmodel);

  % ---- PHASE 2: INFERENCE (Right Column) ----
  \node[label, text=cvprblue] (inflabel) at (5.5, 2.0) {Phase 2: Inference (per test image)};

  \node[draw=blue!30, dashed, rounded corners=8pt, fill=blue!2,
        minimum width=4.8cm, minimum height=7.2cm] (infbox) at (5.5, -1.8) {};

  % Test image
  \node[block, fill=orange!10] (testimg) at (5.5, 1.0) {Test Image\\$I_{\text{test}}$};

  % DINOv2
  \node[block, fill=yellow!10] (dinov2_inf) at (5.5, -0.5) {Frozen\\DINOv2-G};

  % Layer aggregation
  \node[smallbox, fill=yellow!5] (layers_inf) at (5.5, -1.6) {Layers 22--28\\Mean Pooling};

  % Test features
  \node[block, fill=green!8] (testfeats) at (5.5, -2.7) {Patch Feature Map\\$\mathbf{X}_{\text{test}}$};

  % MSAE Reconstruction
  \node[modelblock, fill=cvprblue!12, minimum width=2.8cm, minimum height=1.2cm] (msae_inf) at (5.5, -4.5) {\textbf{MSAE Reconstruction}\\[2pt]
    \footnotesize$\hat{\mathbf{x}}_p = G^*(\text{TopK}_{k_m}(\mathbf{z}_p))$};

  % Multi-scale residual
  \node[block, fill=red!8, minimum width=2.8cm] (residual) at (5.5, -5.8) {Multi-Scale Residual\\[1pt]
    \footnotesize$S(\mathbf{x}_p) = \|\mathbf{x}_p - \hat{\mathbf{x}}_p\|_2^2$};

  % Output: Anomaly map
  \node[block, fill=red!15, minimum width=2.4cm] (anomap) at (5.5, -6.8) {Anomaly Map\\$M \in \mathbb{R}^{h \times w}$};

  % Arrows for inference phase
  \draw[arr] (testimg) -- (dinov2_inf);
  \draw[arr] (dinov2_inf) -- (layers_inf);
  \draw[arr] (layers_inf) -- (testfeats);
  \draw[arrblue] (testfeats) -- (msae_inf);
  \draw[arr] (msae_inf) -- (residual);
  \draw[arr] (residual) -- (anomap);

  % ---- CROSS-PHASE ARROW: model transfer ----
  \draw[arrblue, dashed, line width=1.2pt] (savedmodel.east) -- ++(1.1, 0) |- (msae_inf.west);
  \node[note, text=cvprblue] at (2.8, -5.2) {load};

  % ---- OUTPUT EXPANSION (right side of inference) ----
  % Image-level
  \node[bigblock, fill=red!20, minimum width=2.1cm] (imgscore) at (9.2, -5.0) {Image-Level Score\\[2pt]
    \footnotesize$\displaystyle s_{\text{img}} = \text{TVaR}(M)$};

  % Pixel-level
  \node[bigblock, fill=red!20, minimum width=2.1cm] (pixscore) at (9.2, -6.8) {Pixel-Level Mask\\[2pt]
    \footnotesize Upsample + Smooth};

  \draw[arr] (anomap.east) -- ++(0.8, 0) |- (imgscore.west);
  \draw[arr] (anomap.east) -- ++(0.8, 0) |- (pixscore.west);

  % ---- ANNOTATIONS ----
  % MSAE detail callout
  \node[note, text=black!50] at (0, -8.0) {};

  % Legend-like note
  \node[note] at (2.75, -8.3) {\textbf{MSAE variants:} UW = uniform weighting, RW = reverse weighting};

  % ---- DECORATIVE LABELS ----
  \node[font=\footnotesize\sffamily, rotate=90, text=blue!40] at (-2.8, -2.9) {Training Stage};
  \node[font=\footnotesize\sffamily, rotate=90, text=blue!40] at (8.3, -2.9) {Testing Stage};

  \end{tikzpicture}
  \caption{\textbf{Overview of the MSAE-AD framework.}
  \textbf{Left (Fitting):} $k$ normal images are passed through a frozen DINOv2-G backbone; features from layers 22--28 are mean-pooled to form normal patch features $\mathcal{X}_{\text{normal}}$.
  An MSAE is trained on these features with $M$ nested granularity levels $\{k_1,\ldots,k_M\}$ using either uniform (UW) or reverse (RW) weighting.
  \textbf{Right (Inference):} A test image follows the same feature extraction pipeline. The trained MSAE reconstructs each patch feature; the multi-scale $\ell_2$ reconstruction residual forms the anomaly map, which is aggregated into an image-level score (via TVaR) and upsampled into a pixel-level localization mask.}
  \label{fig:framework}
\end{figure*}
